\chapter{Overview}
\label{chap:overview}

This blueprint is organized by theorem ownership and proof dependency for Ji--Natarajan--Vidick--Wright--Yuen, ``Quantum soundness of the classical low individual degree test''~\cite{Ji2020LowIndividualDegree}, rather than by the raw TeX input order. The formal target is Theorem~\ref{thm:main-formal}. The proof closes by induction on the ambient dimension $m$. The quantitative input to that induction is split into four blocks: the comparison lemmas of Chapter~\ref{chap:preliminaries}, the projectivization tools of Chapter~\ref{chap:projective}, the variance estimate of Chapters~\ref{chap:expansion} and~\ref{chap:variance}, and the self-improvement, commutativity, and pasting chain of Chapters~\ref{chap:self-improvement}, \ref{chap:commutativity}, and~\ref{chap:pasting}.

\begin{theorem}[Raz--Safra]\label{thm:raz-safra}
  Suppose Provers~$\mathrm{A}$ and~$\mathrm{B}$ pass the $k=2$ surface-versus-point low-degree test with probability $1-\eps$. Then there exists a degree-$d$ polynomial $g \colon \F_q^m \to \F_q$ such that
  \[
    \Prb_{u \sim \F_q^m}[g(u)=a] \ge 1-\eps-\poly(m)\cdot\poly(d/q).
  \]
\end{theorem}
\begin{proof}
  See Raz--Safra [RS97].
\end{proof}

\begin{theorem}[Polishchuk--Spielman]\label{thm:classical-test-soundness}
  Suppose Provers~$\mathrm{A}$ and~$\mathrm{B}$ pass the low individual degree test with probability $1-\eps$. Then there exists a polynomial $g \colon \F_q^m \to \F_q$ with individual degree $d$ such that
  \[
    \Prb_{u \sim \F_q^m}[g(u)=a] \ge 1-\poly(m)\cdot\bigl(\poly(\eps)+\poly(d/q)\bigr).
  \]
\end{theorem}
\begin{proof}
  See Polishchuk--Spielman [PS94].
\end{proof}

\begin{theorem}[Main theorem, informal]\label{thm:main-informal}
  If a two-prover projective strategy passes the $(m,q,d)$ low individual degree test with probability $1-\eps$, then there are global polynomial measurements whose evaluations agree with the point answers except with error $\poly(m)(\poly(\eps)+\poly(d/q))$.
\end{theorem}
\begin{proof}
  \uses{thm:main-formal}
  This is the informal form of Theorem~\ref{thm:main-formal}.
\end{proof}
